\documentclass[10pt]{article}

\usepackage[utf8]{inputenc}

\usepackage[T1]{fontenc}

\usepackage{amsmath}

\usepackage{amsfonts}

\usepackage{amssymb}

\usepackage[version=4]{mhchem}

\usepackage{stmaryrd}

\usepackage{graphicx}

\usepackage[export]{adjustbox}

\graphicspath{ {./images/} }



\begin{document}

пересекающиеся определенные над $k$ максимальныс вполне изотропные (относительно симметрической билинейной формы в $V$, ассоциированной с квадратичной формой $Q)$ подпространства в $V, C_{L}$ - подалгебра в алгебре Клиффорда $C=C(Q)$, порожденная подпространством $L \subset V$, и $e_{M} \in C$ - произведение $m$ векторов, составляющих определенный над $k$ базис пространства $M$. Если $n=2 m$ четно, то С. п. реализуется в простом левом идеале $C e_{M}$ и действует там с помощью левых сдвигов: $\rho(s) x=s x\left(s \in C^{+}, x \in C e_{M}\right)$. Далее соответствие $x \mapsto x e_{M}$ определяет изоморфизм векторных пространств $C_{L} \rightarrow C e_{M}$, что позволяет реализовать С. $C_{L}$, естественно изоморфном внепней алгебре над пространством $L$. При этом полуспинорные представления $\rho^{\prime}$ и $\rho^{\prime \prime}$ реализуются в инвариантных $2^{m-1}$-мерных подпространствах $C_{L} \cap C^{+}$и $C_{L} \cap C^{-}$.



Если $n$ нечетно, то пространство $V$ можно включить в $(n+1)$-мерное вскторное пространство $V_{1}=V \oplus k \varepsilon$ над $k$ и определить в $V_{1}$ квадратичную форму $Q_{1}$, положив $Q_{1}(v+\lambda \varepsilon)=Q(v)-\lambda^{2}$ при всех $v \in V$ и $\lambda \in k$. При этом $Q_{1}$ - определенная над $k$ невырожденная квадратичная форма максимального индекса Витта на четномерном векторном пространстве $V_{1}$. С. п. алгебры $C+(Q)$ (группы $\left.\operatorname{Spin}_{n}(Q)\right)$ получается путем ограничения любого из полуспинорных представлений алгебры $C+\left(Q_{1}\right)$ (группы $\left.\operatorname{Spin}_{n+1}\left(Q_{1}\right)\right)$ на подалгебру $C+(Q)$ (соответственно подгруппу $\operatorname{Spin}_{n}(Q)$ ).



В случае, когда $3 \leqslant n \leqslant 14, \quad$ а $k-$ алгөбраически замкнутое поле характеристики 0, решена задача классификации спиноров (см. [4], [8], [9]), к-рая состоит в 1) описании всех орбит группы $\rho\left(\operatorname{Spin}_{n}\right)$ в пространстве спиноров, т. е. указании в каждої орбите нек-рого единственного представителя, 2) вычислении стабилизаторов групиы $\operatorname{Spin}_{n}$ в каждом из этих представителей, 3) описании алгебры инвариантов линейной группы $\rho\left(\operatorname{Spin}_{n}\right)$.



Существование спинорных и полуспинорных представлений алгеб́р Ли $\mathfrak{s v}_{n}$ групп $\operatorname{Spin}_{n}$ было открыто Э. Картаном (É. Cartan) в 1913, когда он классифицировал все неприводимые конечномерные представленил простых алгебр Ли [6]. Впоследствии, в 1935 P. Брауэр (R. Brauer) и Г. Вейль (H. Weyl) описали с пинорные и полуспинорные представления в терминах алгеб́р Клиффорда [5]. П. Дирак (Р. Dirac, [3]) обнаружил, гто при помопц сппноров в квантовой механике огисывается вращение электрона.



.Tum.: [1] Б у р б а к и Н., Алгебра. Модули, кольца, формы, пер. с франц., М., 1966; [2]' В е й л ь Г., Классические группы, их инварианты и представления, пер. с нем., М., 1947; пзд., М., 1979; '[4] II о П о в B. Ј., «Тр. Моск. матем. об-ва», 1978, т. 37, c. 173-217; [5] B r a u e r R., W e y l H., "Amer. J. Math.", 1935, v. $57, \mathrm{~N} 2$, p. $425-49$; [6] C a rta n E., CEuvres complètes, v. 1, pt 1, P., 1952 , p. $355-98 ;[7]$ C h e v a l1 c y C., The algebraic theory of spinors, N. Y., $1954 ;$ [8] G a t-

t i V., V i n i b e r g h 1 E., "Advances in Math.", 1978, v. 30, t i V., V in i b e r g h i E., "Advances in Math.», 1978, v. 30,

№ 2, p. $137-55 ;[9]$ I g u a J.-i., "Amer. J. Math.», 1970, v. 92 , № 4, p. 997-1028. B. \#. IIonos.



СПИРАЛИ - плоские кривые, к-рые обычно обходят вокруг одної (или нескольких точек), приближаясь или удаляясь от нее.



Среди С. выделяют алгебраич. С. и псевдосширали.



А лге браически е с пи р али-спирали, уравнения к-рых в полярных координатах являются алгебраическими относительно переменных $\rho$ и $\varphi$. Ii алгебраическим С. относятся: гиперболическая спираль, архимедова спираль, Галилея спираль, Ферма спираль, параболическая спираль, жезл.



П с в в о с пи р а ли-спираль, натуральные уравнения к-рых могут быть записаны в виде



\[

r=a s^{m}

\]



где $r$ - радиус кривизны, $s$ - длина дуги. Пріь $m=1$ псевдоспираль наз. логариялической спиралью, при $m=-1$ - Корню спиралью, при $m=1 / 2$ является әвольвентой окружности.



Лит.: [1] С а в е ло в А. А., Плоские кривые, М., 1960. Д. Д. Соколов.



SICI-CIИРАЛЬ - плоская кривая, уравнение к-рой в декартовых прямоугольных координатах $(x, y)$ имеет вид



\[

x=\operatorname{ci}(t), y=\operatorname{si}(t),

\]



где ci - интегральный косинус, si - интегральный синус, $t$ - действительный параметр (рис.). Дли-



\begin{center}

\includegraphics[max width=\textwidth]{2023_07_09_7b904f9ff6dbb3a8bff4g-1}

\end{center}



на дуги от точки $t=0$ до точки $t_{0}$ равна $\ln t_{0}$, а кривизна $x=t_{0}$.



Лит.: [1] Я н к е Е., Э м д е Ф., Л ё шा Ф., Специальные функции. Формулы, графики, табллицы, пер. с нем., 2 изд., М.,

Д. Д. Соколов.



СПИРМЕНА КОЭФФИЦИЕНТ РАНГОВОЙ КОРРЕЛЯЩИИ - мера зависимости двух случайных величин (признаков) $X$ и $Y$, основанная на ранжировании независимых результатов наблюдений $\left(X_{1}, Y_{1}\right), \ldots,\left(X_{n}, Y_{n}\right)$. Если ранги значений $X$ расположены в естественном порядке $i=1, \ldots, n$, а $R_{i}-$ ранг $Y$, соответствующий той паре $(X, Y)$, для к-роӥ ранг $X$ равен $i$, то С. к. р. к. определяется क्रормулой



\[

r_{s}=\frac{12}{n\left(n^{2}-1\right)} \sum_{i=1}^{n}\left(i-\frac{n+1}{2}\right)\left(R_{i}-\frac{n+1}{2}\right)

\]



или, что равносильно,



\[

r_{s}=1-\frac{6 \sum_{i=1}^{n} d_{i}^{2}}{n\left(n^{2}-1\right)},

\]



где $d_{i}$ - разность между рангами $X_{i}$ и $Y_{i}$. Значение $r_{s}$ меняется от -1 до +1 , причем $r_{s}=+1$, когда последовательности рангов полностью совпадают, т. е. $i=$ $=R_{i}, i=1, \ldots, n$, и $r_{s}=-1$, когда последовательности рангов полностью противоположны, т. е. $i=(n+1)-R_{i}$, $i=1, \ldots, n$. С. к. Г. к., как и любая другая ранговая статистика, применяется для проверки гипотезы независимости двух признаков. Еслии признаки независимы, то $\mathrm{E} r_{s}=0, \mathrm{D} r_{s}=\frac{1}{n-1}$. Таким образом, но величине отклонения $r_{s}$ от нуля монко сделать вывод о зависимости или незанисимости признаков. Для построеиия соответствующего критерия вычисляется распределение $r_{s}$ для независимых признаков $X$ и $Y$. При $4 \leqslant n \leqslant$ $\leqslant 10$ используют таблицы точного распределения (см. [2], [4]), а при $n>10$ можно воспользоваться, напр., тем, что случайная величина $\sqrt{n-1} r_{s}$ при $n \rightarrow \infty$ распределена асимптотически нормально с параметрами $(0,1)$. В последнем случае гипотеза независимости отвергается, если $\left|r_{s}\right|>u_{1-\frac{\alpha}{2}} / \sqrt{n-1}$, где $u_{1-\frac{\alpha}{2}}$ есть корень уравнения $₫(u)=1-\frac{\alpha}{2}(\Phi(u)$ - функция стандартного нормального распределения).





\end{document}